% ==============================================================================
% 页面布局与基础设置 (Page Layout & Core Settings)
% ==============================================================================
\usepackage{geometry}
\geometry{
    paper=a3paper,
    landscape,
    left=0.3cm, right=0.3cm, top=0.3cm, bottom=0.3cm
}

\usepackage[fontsize=10pt]{fontsize} % 全局 9pt 字号
\linespread{1.1}                   % 紧凑行距
\setlength{\parindent}{0pt}         % 禁止全局首行缩进

\usepackage{multicol}               % 多栏排版
\setlength{\columnsep}{6pt}         % 栏间距
\setlength{\columnseprule}{0pt}     % 栏分割线

% ==============================================================================
% 语言、字体与数学支持 (Fonts & Mathematics)
% ==============================================================================
\usepackage[scheme=plain]{ctex}     % 中文支持
\usepackage{fontspec}               % 字体管理
\IfFontExistsTF{Times New Roman}    % 西文字体；没装 Times New Roman 时用 TeX Gyre Termes（Times 克隆，TeX Live 自带）
  {\setmainfont{Times New Roman}}
  {\setmainfont{texgyretermes}[Extension=.otf, UprightFont=*-regular, BoldFont=*-bold,
                               ItalicFont=*-italic, BoldItalicFont=*-bolditalic]}

\usepackage{amsmath,amssymb}        % 基础数学公式
\usepackage{mathtools}              % 扩展数学增强
\usepackage{calc}                   % 长度计算宏包

% 常用数学简写定义
\newcommand{\E}{\mathbb{E}}
\newcommand{\Var}{\mathrm{Var}}
\newcommand{\Cov}{\mathrm{Cov}}
\newcommand{\bs}[1]{\boldsymbol{#1}}

% ==============================================================================
% 样式定制 (Styling: Sections, Lists, Colors)
% ==============================================================================
\usepackage{xcolor}
\definecolor{myblue}{RGB}{0, 0, 165}
\definecolor{blue}{RGB}{0, 133, 255}
\definecolor{red}{RGB}{210, 0, 0}
\definecolor{grey}{RGB}{127, 127, 127}

% --- 列表样式定制 ---
\usepackage{enumitem}
\setlist[itemize]{leftmargin=*, noitemsep, topsep=0pt, parsep=0pt} % 极紧凑列表

% --- 标题样式定制 (Titlesec) ---
\usepackage[explicit]{titlesec}

% 重定义 \section (带下划线风格)
\titleformat{\section}
  {\normalfont\large\bfseries\fontsize{12pt}{12pt}\color{myblue}}
  {\thesection.}
  {2pt}
  {#1}
  [{\color{myblue}\titlerule[0.5pt]}]
\titlespacing*{\section}{0pt}{2pt}{2pt}

% 重定义 \subsection (蓝色加粗风格)
\titleformat{\subsection}
  {\normalfont\bfseries\fontsize{10pt}{10pt}\color{myblue}} 
  {\thesubsection.}
  {2pt}
  {#1}
\titlespacing*{\subsection}{0pt}{1pt}{1pt}

% 设置 subsubsection 的格式
\titleformat{\subsubsection}
  {\normalfont\normalsize\bfseries\color{myblue}} % 格式：正常字体、正常大小、加粗
  {\thesubsubsection}               % 标签：显示章节编号
  {2pt}                             % 编号与标题之间的距离
  {#1}                                % 标题前的额外代码
\titlespacing*{\subsubsection}
  {0pt}      % 左边缩进
  {0pt}      % 段前间距 (设置为0或很小的值)
  {0pt}      % 段后间距 (设置为0，紧贴正文)

\titlespacing*{\paragraph}{0pt}{0pt}{4pt}

\titlespacing*{\subparagraph}{0pt}{0pt}{4pt}

\usepackage{enumitem}
\setlist{noitemsep, topsep=0pt} %调整列表格式
\usepackage{caption}
% ==============================================================================
% 代码块配置 (Code Listings: Stata Support)
% ==============================================================================
\usepackage{listings}

\lstset{
  basicstyle=\ttfamily\footnotesize,
  tabsize=4,
  columns=fullflexible,
  breaklines=true,
  breakatwhitespace=true,
  frame=none,
  % --- 空格处理 ---
  showspaces=false,               % 全局空格不显示特殊符号
  showstringspaces=false,         % 字符串内的空格不显示特殊符号
  showtabs=false,                 % 制表符不显示特殊符号
  % --- 行号设置 ---
  numbers=left, 
  numberstyle=\tiny\color{gray}, 
  numbersep=4pt, 
  xleftmargin=5pt,
  keywordstyle=\color{blue},
  commentstyle=\color{gray},
  stringstyle=\color{red},
}

\lstdefinelanguage{Stata}{
  morekeywords={
    regress, reg, gen, generate, replace, drop, keep,
    summarize, sum, tabulate, tab, egen,
    if, in, using, clear,
    by, bysort, foreach, forvalues,
    do, run, log, use, save, merge, append,
    xtset, xtreg, xtivreg,
    vce, cluster, robust
  },
  sensitive=true,
  morecomment=[l]{*},
  morestring=[b]",
}

% ==============================================================================
% 文档全局间距微调 (Math Spacing Tuning)
% ==============================================================================
% ==============================================================================
% Cheatsheet 追加：表格、行内代码、强调、图
% ==============================================================================
\usepackage{tabularx,booktabs,array}
\newcolumntype{Y}{>{\raggedright\arraybackslash}X}  % 自动换行的左对齐列
\setlength{\tabcolsep}{2pt}
\setlength{\aboverulesep}{0.4pt}
\setlength{\belowrulesep}{0.4pt}
\setlist[enumerate]{leftmargin=*, noitemsep, topsep=0pt, parsep=0pt}
\setlist[description]{leftmargin=0pt, noitemsep, topsep=0pt, parsep=0pt}

\newcommand{\cd}[1]{\texttt{\detokenize{#1}}}       % 行内代码（可含 _ 和 ->）
\newcommand{\imp}[1]{\textbf{\textcolor{red}{#1}}}   % 重点
\newcommand{\ez}{\textbf{\textcolor{red}{【易错】}}}
\newcommand{\bk}{\textbf{\textcolor{red}{【必考】}}}
\newcommand{\Q}[1]{\par\textbf{\textcolor{myblue}{#1}}\par}  % 问题小标题

\lstset{numbers=none, xleftmargin=0pt, basicstyle=\ttfamily\footnotesize,
        aboveskip=1pt, belowskip=1pt}

\usepackage{tikz}
\usetikzlibrary{arrows.meta,positioning,calc}

\AtBeginDocument{
    \setlength{\abovedisplayskip}{1pt}      % 公式上方间距
    \setlength{\belowdisplayskip}{1pt}      % 公式下方间距
    \setlength{\abovedisplayshortskip}{0pt} % 短公式上方间距
    \setlength{\belowdisplayshortskip}{0pt} % 短公式下方间距
    \setlength{\jot}{0pt}                   % 多行公式行间距
}
% --- 可跨栏续排的表格 ---
% 每一行排成一个独立的单行表格，所有行使用相同的固定列宽，因此上下对齐；
% 行与行之间允许分栏，放不下时剩下的行接到下一栏。
\usepackage{needspace}  % 章标题前预留空间，避免标题孤悬栏底
\newcolumntype{P}[1]{>{\raggedright\arraybackslash}p{\dimexpr #1\linewidth-2\tabcolsep\relax}}
\newcommand{\tspec}{}
% \begin{btab}{列格式} ... \end{btab}；表头用 \thead{...}，每行用 \trow{...}
\newenvironment{btab}[1]
  {\par\needspace{4\baselineskip}\footnotesize\renewcommand{\tspec}{#1}\setlength{\parskip}{0pt}\vspace{2pt}%
   \noindent\rule{\linewidth}{0.8pt}\par\nointerlineskip}
  {\par\nointerlineskip\noindent\rule{\linewidth}{0.8pt}\par\vspace{2pt}}
\newcommand{\btabular}[2]{\begin{tabular}{@{}#1@{}}\rule{0pt}{2.6ex}#2\end{tabular}}
\newcommand{\trow}[1]{\par\nointerlineskip\noindent
  \expandafter\btabular\expandafter{\tspec}{#1}\par}
\newcommand{\thead}[1]{\trow{#1}\nointerlineskip\noindent\rule{\linewidth}{0.4pt}\par\nopagebreak}
